\documentclass[10pt,a4paper]{amsart}
\usepackage{amsmath,amssymb,amsthm,graphicx}
\usepackage{verbatim}
\newenvironment{reference}{\comment}{\endcomment}
%\newenvironment{reference}{}{}
\newenvironment{slogan}{\comment}{\endcomment}
\newenvironment{history}{\comment}{\endcomment}
% For commutative diagrams we use Xy-pic
\usepackage[all]{xy}

% We use 2cell for 2-commutative diagrams.
\xyoption{2cell}
% Theorem environments.
%
\theoremstyle{plain}
\newtheorem{theorem}[subsection]{Theorem}
\newtheorem{proposition}[subsection]{Proposition}
\newtheorem{lemma}[subsection]{Lemma}

\theoremstyle{definition}
\newtheorem{definition}[subsection]{Definition}
\newtheorem{example}[subsection]{Example}
\newtheorem{exercise}[subsection]{Exercise}
\newtheorem{situation}[subsection]{Situation}

\theoremstyle{remark}
\newtheorem{remark}[subsection]{Remark}
\newtheorem{remarks}[subsection]{Remarks}

\numberwithin{equation}{subsection}

% Macros
%
\def\lim{\mathop{\mathrm{lim}}\nolimits}
\def\colim{\mathop{\mathrm{colim}}\nolimits}
\def\Spec{\mathop{\mathrm{Spec}}}
\def\Hom{\mathop{\mathrm{Hom}}\nolimits}
\def\Ext{\mathop{\mathrm{Ext}}\nolimits}
\def\SheafHom{\mathop{\mathcal{H}\!\mathit{om}}\nolimits}
\def\SheafExt{\mathop{\mathcal{E}\!\mathit{xt}}\nolimits}
\def\Sch{\mathit{Sch}}
\def\Mor{\mathop{\mathrm{Mor}}\nolimits}
\def\Ob{\mathop{\mathrm{Ob}}\nolimits}
\def\Sh{\mathop{\mathit{Sh}}\nolimits}
\def\NL{\mathop{N\!L}\nolimits}
\def\CH{\mathop{\mathrm{CH}}\nolimits}
\def\proetale{{pro\text{-}\acute{e}tale}}
\def\etale{{\acute{e}tale}}
\def\QCoh{\mathit{QCoh}}
\def\Ker{\mathop{\mathrm{Ker}}}
\def\Im{\mathop{\mathrm{Im}}}
\def\Coker{\mathop{\mathrm{Coker}}}
\def\Coim{\mathop{\mathrm{Coim}}}

% Boxtimes
%
\DeclareMathSymbol{\boxtimes}{\mathbin}{AMSa}{"02}

%
% Macros for moduli stacks/spaces
%
\def\QCohstack{\mathcal{QC}\!\mathit{oh}}
\def\Cohstack{\mathcal{C}\!\mathit{oh}}
\def\Spacesstack{\mathcal{S}\!\mathit{paces}}
\def\Quotfunctor{\mathrm{Quot}}
\def\Hilbfunctor{\mathrm{Hilb}}
\def\Curvesstack{\mathcal{C}\!\mathit{urves}}
\def\Polarizedstack{\mathcal{P}\!\mathit{olarized}}
\def\Complexesstack{\mathcal{C}\!\mathit{omplexes}}
% \Pic is the operator that assigns to X its picard group, usage \Pic(X)
% \Picardstack_{X/B} denotes the Picard stack of X over B
% \Picardfunctor_{X/B} denotes the Picard functor of X over B
\def\Pic{\mathop{\mathrm{Pic}}\nolimits}
\def\Picardstack{\mathcal{P}\!\mathit{ic}}
\def\Picardfunctor{\mathrm{Pic}}
\def\Deformationcategory{\mathcal{D}\!\mathit{ef}}
\usepackage{cleveref}
\crefname{originallemma}{original lemma}{original lemmas}
\Crefname{originallemma}{Original lemma}{Original lemmas}
\makeatletter
\newlabel{lemma-stalk-tensor-product}{{lemma-stalk-tensor-product}{source}{}{sourcepointer-1}{}}
\newlabel{lemma-stalk-tensor-product@cref}{{[originallemma][0][]lemma-stalk-tensor-product}{[1][0][]source}{}{}{}}
\makeatother
\title{Flatness of a module on a ringed space is detected on the stalks}
\author{The Stacks Project Authors}
\date{stacks-project:modules.tex:local-snapshot-20261012; GNU FDL 1.2 or later}
\begin{document}
\maketitle
\noindent\textit{Editorial note.} All delivered mathematical text is the author's own, copied from the pinned Stacks Project chapter source sources/stacks-modules/pinned-original-source.tex; no statement or proof step was rewritten, reordered or removed.\par\medskip
\noindent\textit{Editorial note.} Preamble provenance: the Stacks Project keeps one single shared preamble for every chapter; this card takes it directly from the staged sources/stacks-modules/preamble.tex as three verbatim spans (lines 8-12, 14-18 and 154-217), so no derived file was written and nothing was added or redefined. The wrapper class and the base packages are supplied by the assembler.\par\medskip
\noindent\textit{Editorial note.} Bibliography provenance: neither the counted statement nor its proof contains a citation (this lane rescans the whole span before writing the card), so the card carries no bibliography and no entry was reconstructed.\par\medskip
\noindent\textit{Editorial note.} Selection provenance: this card was opened under the file-based control contract (CONTROL.md), with the mechanical pre-card gate run on its statement and proof interval before the card was written.\par\medskip
\noindent\textit{Allowed standard prerequisites (editorial).}\begin{itemize}
\item The description of the stalk of a tensor product of modules as the tensor product of the stalks. Invoked twice in the counted proof (source lines 2781-2783 and 2791-2792) as a pointer to a lemma of this chapter. Transfers exactness of a sequence of modules to exactness after tensoring, stalk by stalk.
\item The skyscraper sheaf at a point associated to a module over the local ring. Used inside the counted proof (source lines 2784-2795). Realises an arbitrary module over the local ring as the stalk of a sheaf, which tests flatness of the stalk.
\end{itemize}\medskip
\ifdefined\hypertarget\hypertarget{sourcepointer-1}{}\fi
\noindent\textit{Source reference.} Original lemma lemma-stalk-tensor-product: sources/stacks-modules/pinned-original-source.tex, source lines 2604--2604. Navigation pointer only: the counted text refers to the Stacks item lemma-stalk-tensor-product, which belongs to another chapter of the project and is not part of this card. The pointer carries the target's own label as its number because this chapter's staged snapshot does not contain the other chapter's tag table; its content is not reproduced here.\par\medskip
\begin{lemma}
\label{lemma-flat-stalks-flat}
Let $(X, \mathcal{O}_X)$ be a ringed space.
An $\mathcal{O}_X$-module $\mathcal{F}$ is flat
if and only if the stalk $\mathcal{F}_x$ is a flat
$\mathcal{O}_{X, x}$-module for all $x \in X$.
\end{lemma}

\begin{proof}
Assume $\mathcal{F}_x$ is a flat $\mathcal{O}_{X, x}$-module for all
$x \in X$. In this case, if $\mathcal{G} \to \mathcal{H} \to \mathcal{K}$
is exact, then also
$\mathcal{G} \otimes_{\mathcal{O}_X} \mathcal{F} \to
\mathcal{H} \otimes_{\mathcal{O}_X} \mathcal{F} \to
\mathcal{K} \otimes_{\mathcal{O}_X} \mathcal{F}$
is exact because we can check exactness at stalks and because
tensor product commutes with taking stalks, see
Lemma \ref{lemma-stalk-tensor-product}.
Conversely, suppose that $\mathcal{F}$ is flat, and let $x \in X$.
Consider the skyscraper sheaves $i_{x, *} M$ where $M$ is a
$\mathcal{O}_{X, x}$-module. Note that
$$
M \otimes_{\mathcal{O}_{X, x}} \mathcal{F}_x =
\left(i_{x, *} M \otimes_{\mathcal{O}_X} \mathcal{F}\right)_x
$$
again by
Lemma \ref{lemma-stalk-tensor-product}.
Since $i_{x, *}$ is exact, we see that the fact that $\mathcal{F}$
is flat implies that $M \mapsto M \otimes_{\mathcal{O}_{X, x}} \mathcal{F}_x$
is exact. Hence $\mathcal{F}_x$ is a flat $\mathcal{O}_{X, x}$-module.
\end{proof}

\end{document}
